% Options for packages loaded elsewhere
\PassOptionsToPackage{unicode}{hyperref}
\PassOptionsToPackage{hyphens}{url}
\PassOptionsToPackage{dvipsnames,svgnames,x11names}{xcolor}
%
\documentclass[
  letterpaper,
  DIV=11,
  numbers=noendperiod]{scrartcl}

\usepackage{amsmath,amssymb}
\usepackage{iftex}
\ifPDFTeX
  \usepackage[T1]{fontenc}
  \usepackage[utf8]{inputenc}
  \usepackage{textcomp} % provide euro and other symbols
\else % if luatex or xetex
  \usepackage{unicode-math}
  \defaultfontfeatures{Scale=MatchLowercase}
  \defaultfontfeatures[\rmfamily]{Ligatures=TeX,Scale=1}
\fi
\usepackage{lmodern}
\ifPDFTeX\else  
    % xetex/luatex font selection
\fi
% Use upquote if available, for straight quotes in verbatim environments
\IfFileExists{upquote.sty}{\usepackage{upquote}}{}
\IfFileExists{microtype.sty}{% use microtype if available
  \usepackage[]{microtype}
  \UseMicrotypeSet[protrusion]{basicmath} % disable protrusion for tt fonts
}{}
\makeatletter
\@ifundefined{KOMAClassName}{% if non-KOMA class
  \IfFileExists{parskip.sty}{%
    \usepackage{parskip}
  }{% else
    \setlength{\parindent}{0pt}
    \setlength{\parskip}{6pt plus 2pt minus 1pt}}
}{% if KOMA class
  \KOMAoptions{parskip=half}}
\makeatother
\usepackage{xcolor}
\usepackage[margin=1in]{geometry}
\setlength{\emergencystretch}{3em} % prevent overfull lines
\setcounter{secnumdepth}{5}
% Make \paragraph and \subparagraph free-standing
\ifx\paragraph\undefined\else
  \let\oldparagraph\paragraph
  \renewcommand{\paragraph}[1]{\oldparagraph{#1}\mbox{}}
\fi
\ifx\subparagraph\undefined\else
  \let\oldsubparagraph\subparagraph
  \renewcommand{\subparagraph}[1]{\oldsubparagraph{#1}\mbox{}}
\fi


\providecommand{\tightlist}{%
  \setlength{\itemsep}{0pt}\setlength{\parskip}{0pt}}\usepackage{longtable,booktabs,array}
\usepackage{calc} % for calculating minipage widths
% Correct order of tables after \paragraph or \subparagraph
\usepackage{etoolbox}
\makeatletter
\patchcmd\longtable{\par}{\if@noskipsec\mbox{}\fi\par}{}{}
\makeatother
% Allow footnotes in longtable head/foot
\IfFileExists{footnotehyper.sty}{\usepackage{footnotehyper}}{\usepackage{footnote}}
\makesavenoteenv{longtable}
\usepackage{graphicx}
\makeatletter
\def\maxwidth{\ifdim\Gin@nat@width>\linewidth\linewidth\else\Gin@nat@width\fi}
\def\maxheight{\ifdim\Gin@nat@height>\textheight\textheight\else\Gin@nat@height\fi}
\makeatother
% Scale images if necessary, so that they will not overflow the page
% margins by default, and it is still possible to overwrite the defaults
% using explicit options in \includegraphics[width, height, ...]{}
\setkeys{Gin}{width=\maxwidth,height=\maxheight,keepaspectratio}
% Set default figure placement to htbp
\makeatletter
\def\fps@figure{htbp}
\makeatother

\KOMAoption{captions}{tableheading}
\makeatletter
\@ifpackageloaded{caption}{}{\usepackage{caption}}
\AtBeginDocument{%
\ifdefined\contentsname
  \renewcommand*\contentsname{Table of contents}
\else
  \newcommand\contentsname{Table of contents}
\fi
\ifdefined\listfigurename
  \renewcommand*\listfigurename{List of Figures}
\else
  \newcommand\listfigurename{List of Figures}
\fi
\ifdefined\listtablename
  \renewcommand*\listtablename{List of Tables}
\else
  \newcommand\listtablename{List of Tables}
\fi
\ifdefined\figurename
  \renewcommand*\figurename{Figure}
\else
  \newcommand\figurename{Figure}
\fi
\ifdefined\tablename
  \renewcommand*\tablename{Table}
\else
  \newcommand\tablename{Table}
\fi
}
\@ifpackageloaded{float}{}{\usepackage{float}}
\floatstyle{ruled}
\@ifundefined{c@chapter}{\newfloat{codelisting}{h}{lop}}{\newfloat{codelisting}{h}{lop}[chapter]}
\floatname{codelisting}{Listing}
\newcommand*\listoflistings{\listof{codelisting}{List of Listings}}
\makeatother
\makeatletter
\makeatother
\makeatletter
\@ifpackageloaded{caption}{}{\usepackage{caption}}
\@ifpackageloaded{subcaption}{}{\usepackage{subcaption}}
\makeatother
\ifLuaTeX
  \usepackage{selnolig}  % disable illegal ligatures
\fi
\usepackage{bookmark}

\IfFileExists{xurl.sty}{\usepackage{xurl}}{} % add URL line breaks if available
\urlstyle{same} % disable monospaced font for URLs
\hypersetup{
  pdftitle={Informe Técnico: Sistema de Análisis e Inteligencia de Negocios - FIFA},
  colorlinks=true,
  linkcolor={blue},
  filecolor={Maroon},
  citecolor={Blue},
  urlcolor={Blue},
  pdfcreator={LaTeX via pandoc}}

\title{Informe Técnico: Sistema de Análisis e Inteligencia de Negocios -
FIFA}
\author{}
\date{2026-07-23}

\begin{document}
\maketitle

\renewcommand*\contentsname{Table of contents}
{
\hypersetup{linkcolor=}
\setcounter{tocdepth}{3}
\tableofcontents
}
\section{1. Objetivos del Proyecto}\label{objetivos-del-proyecto}

\subsection{Objetivo General}\label{objetivo-general}

Implementar un sistema de Inteligencia de Negocios basado en un proceso
ETL automatizado para el análisis de datos históricos de la Copa Mundial
de Fútbol, permitiendo la extracción, transformación y carga de
información mediante Python y MySQL, con la finalidad de generar
indicadores clave de rendimiento (KPIs) y visualizarlos mediante un
Dashboard BI interactivo para facilitar el análisis deportivo.

\subsection{Objetivos Específicos}\label{objetivos-especuxedficos}

\begin{itemize}
\item
  Diseñar una arquitectura distribuida utilizando bases de datos y
  servidores independientes para separar la información operativa del
  procesamiento analítico.
\item
  Implementar un proceso ETL mediante Python que permita extraer datos
  históricos de la FIFA, aplicar transformaciones y calcular indicadores
  de rendimiento deportivo.
\item
  Utilizar parámetros de configuración dinámicos para controlar el
  procesamiento de datos y evitar valores definidos directamente dentro
  del código fuente.
\item
  Automatizar la ejecución del flujo de procesamiento mediante
  herramientas de orquestación, garantizando una ejecución organizada y
  repetible del proceso ETL.
\item
  Construir un Dashboard BI interactivo que permita visualizar los KPIs
  generados mediante gráficos especializados y elementos visuales para
  el análisis de resultados históricos.
\end{itemize}

\section{2. Arquitectura y Diseño de la Solución (Valor: 1.5
pts)}\label{arquitectura-y-diseuxf1o-de-la-soluciuxf3n-valor-1.5-pts}

\subsection{Diagrama de Arquitectura}\label{diagrama-de-arquitectura}

\begin{figure}[H]

{\centering \includegraphics[width=0.9\textwidth,height=\textheight]{imagenes/arquitectura.png}

}

\caption{Arquitectura General del Sistema Distribuido}

\end{figure}%

\subsection{Explicación del Flujo
ETL}\label{explicaciuxf3n-del-flujo-etl}

La solución implementada utiliza una arquitectura distribuida compuesta
por servidores Linux Ubuntu, bases de datos MySQL y herramientas de
procesamiento desarrolladas en Python.

El flujo general del sistema se divide en las siguientes etapas:

\begin{enumerate}
\def\labelenumi{\arabic{enumi}.}
\item
  \textbf{Extracción de datos:}

  El proceso ETL inicia mediante la conexión hacia la base de datos que
  contiene la información histórica de la Copa Mundial de Fútbol. Los
  datos son extraídos utilizando Python junto con librerías de conexión
  a bases de datos, permitiendo obtener la información necesaria para el
  análisis.
\item
  \textbf{Transformación de datos:}

  Una vez obtenidos los datos, se aplican procesos de limpieza, filtrado
  y transformación utilizando la librería Pandas. En esta etapa se
  calculan los indicadores clave de rendimiento (KPIs), como
  estadísticas de victorias e índices de rendimiento de las selecciones.
\item
  \textbf{Carga de información analítica:}
\end{enumerate}

\section{3. Explicación Técnico-Analítica (Valor: 1
pt)}\label{explicaciuxf3n-tuxe9cnico-analuxedtica-valor-1-pt}

\subsection{Proceso ETL Implementado}\label{proceso-etl-implementado}

El sistema desarrollado implementa un proceso ETL (Extract, Transform,
Load) encargado de obtener información histórica de la Copa Mundial de
Fútbol, procesarla y almacenarla en una base de datos analítica para su
posterior visualización mediante un Dashboard BI.

La etapa de \textbf{extracción} consiste en la obtención de los datos
históricos almacenados en la base de datos FIFA. Mediante scripts
desarrollados en Python se realiza la conexión hacia la fuente de
información y se recuperan los registros necesarios para el análisis de
rendimiento de las selecciones.

En la etapa de \textbf{transformación}, los datos obtenidos son
procesados utilizando librerías especializadas como Pandas. Durante esta
fase se aplican filtros, limpieza de información y cálculos estadísticos
necesarios para generar indicadores de rendimiento deportivo. Además, se
utilizan parámetros de configuración para controlar dinámicamente las
reglas del procesamiento.

La etapa de \textbf{carga} permite almacenar los resultados procesados
dentro de la base de datos analítica \texttt{bi\_db}, donde se
encuentran las estructuras utilizadas posteriormente por el Dashboard BI
para la generación de visualizaciones interactivas.

\subsection{Control de Ejecuciones e
Idempotencia}\label{control-de-ejecuciones-e-idempotencia}

El proceso ETL fue diseñado considerando el principio de idempotencia,
permitiendo ejecutar nuevamente el flujo de procesamiento sin generar
duplicidad de información ni afectar la consistencia de los datos
almacenados.

Para lograr este comportamiento, la carga de información analítica
mantiene una estructura controlada de actualización de resultados. De
esta manera, cuando el proceso se ejecuta nuevamente bajo los mismos
parámetros de configuración, se obtiene el mismo resultado esperado.

La aplicación de este principio facilita la automatización del proceso
ETL mediante herramientas de orquestación, permitiendo realizar
ejecuciones programadas de manera segura y confiable.

\subsection{Uso de Parámetros de
Configuración}\label{uso-de-paruxe1metros-de-configuraciuxf3n}

El sistema utiliza parámetros de configuración almacenados en la base de
datos de parámetros para evitar valores definidos directamente dentro
del código fuente y permitir mayor flexibilidad en el procesamiento de
información.

Estos parámetros permiten controlar aspectos importantes del análisis,
como:

\begin{itemize}
\item
  \textbf{Año inicial de procesamiento (\texttt{year\_from}):} Permite
  establecer desde qué periodo histórico serán considerados los datos.
\item
  \textbf{Cantidad de registros mostrados (\texttt{top\_n}):} Define el
  número de selecciones que serán incluidas en los rankings principales
  del Dashboard.
\item
  \textbf{Umbrales de evaluación (\texttt{min\_avg\_goals}):} Permite
  establecer valores mínimos utilizados para clasificar características
  de rendimiento ofensivo.
\end{itemize}

El uso de parámetros dinámicos facilita la modificación de reglas de
negocio sin necesidad de cambiar la lógica principal del programa.

\subsection{Cálculo de Indicadores Clave de Rendimiento
(KPIs)}\label{cuxe1lculo-de-indicadores-clave-de-rendimiento-kpis}

Durante la etapa de transformación se generan indicadores que permiten
analizar el desempeño histórico de las selecciones participantes.

Los principales KPIs generados son:

\begin{itemize}
\item
  \textbf{Victorias acumuladas:} Representa la cantidad total de
  partidos ganados por cada selección dentro del periodo analizado.
\item
  \textbf{Índice de rendimiento:} Permite evaluar el desempeño general
  de los equipos combinando los resultados obtenidos durante sus
  participaciones históricas.
\item
  \textbf{Ranking de selecciones:} Organiza los equipos según sus
  indicadores obtenidos, permitiendo identificar las selecciones con
  mejores resultados.
\end{itemize}

Estos indicadores convierten los datos históricos almacenados en
información útil para realizar análisis comparativos mediante el
Dashboard BI.

\subsection{Decisiones Técnicas
Implementadas}\label{decisiones-tuxe9cnicas-implementadas}

Para el desarrollo de la solución se tomaron las siguientes decisiones
técnicas:

\begin{itemize}
\item
  Utilizar \textbf{Python} como lenguaje principal para el desarrollo
  del proceso ETL debido a su capacidad para realizar procesamiento de
  datos, conexión con bases de datos y automatización de tareas.
\item
  Utilizar \textbf{MySQL} como sistema gestor de bases de datos para
  almacenar la información operativa y analítica del proyecto.
\item
  Emplear la librería \textbf{Pandas} para la limpieza, transformación y
  cálculo de métricas debido a su eficiencia en el manejo de datos
  tabulares.
\item
  Implementar un Dashboard BI utilizando \textbf{Python y Plotly},
  permitiendo representar los indicadores mediante gráficos interactivos
  y facilitar la interpretación de resultados.
\item
  Utilizar herramientas de orquestación como \textbf{Apache Airflow}
  para organizar la ejecución del flujo ETL y permitir procesos
  automatizados y repetibles.
\end{itemize}

Estas decisiones permiten construir una solución organizada, escalable y
orientada al análisis de información deportiva histórica.

\section{4. Evidencias de Ejecución (Valor: 1
pt)}\label{evidencias-de-ejecuciuxf3n-valor-1-pt}

En esta sección se presentan las evidencias obtenidas durante la
ejecución del sistema, demostrando el funcionamiento del proceso ETL, la
generación de indicadores y la visualización mediante el Dashboard BI.

\subsection{Ejecución del Proceso
ETL}\label{ejecuciuxf3n-del-proceso-etl}

La siguiente evidencia muestra la ejecución del proceso de extracción,
transformación y carga de datos, donde se observa la conexión con las
fuentes de información, procesamiento de datos y generación de
resultados analíticos.

{[}Captura pendiente: ejecución del proceso ETL{]}

\subsection{Resultados Generados en la Base de Datos
Analítica}\label{resultados-generados-en-la-base-de-datos-analuxedtica}

Se presenta la evidencia de los resultados almacenados en la base de
datos \texttt{bi\_db}, donde se encuentran los indicadores calculados
utilizados posteriormente por el Dashboard BI.

{[}Captura pendiente: tablas o resultados de KPIs generados{]}

\subsection{Dashboard BI Funcionando}\label{dashboard-bi-funcionando}

La siguiente evidencia corresponde al Dashboard BI implementado, donde
se muestran las visualizaciones interactivas generadas a partir de los
KPIs obtenidos durante el proceso ETL.

{[}Captura pendiente: Dashboard BI con gráficos y KPIs{]}

\section{5. Conclusiones}\label{conclusiones}

El desarrollo del sistema de Inteligencia de Negocios permitió
implementar un flujo completo de análisis de datos deportivos mediante
una arquitectura distribuida, integrando bases de datos MySQL,
procesamiento con Python y visualización mediante un Dashboard BI
interactivo.

La implementación del proceso ETL permitió transformar datos históricos
de la Copa Mundial de Fútbol en información analítica mediante la
extracción, procesamiento y carga de indicadores clave de rendimiento.
El uso de parámetros de configuración permitió mejorar la flexibilidad
del sistema, evitando depender de valores definidos directamente dentro
del código.

La aplicación de mecanismos como la idempotencia permitió garantizar
ejecuciones controladas del proceso de procesamiento de datos,
facilitando la automatización y actualización de la información
analítica.

Finalmente, el Dashboard BI desarrollado permitió representar los KPIs
generados mediante gráficos interactivos, facilitando la interpretación
de resultados y proporcionando una herramienta visual para comparar el
rendimiento histórico de las selecciones.



\end{document}
